\makeatletter
\def\input@path{{./}{../}{../../}{../../../}{../../../../}{preamble/}{../preamble/}{../../preamble/}{../../../preamble/}}
\makeatother

\input{preamble_exams.tex}

\title{Decision-Making Under Uncertainty: Exam Solutions}
\author{}
\date{}

\begin{document}
	\maketitle
	\vspace{-2cm}
	\setcounter{tocdepth}{1}
	\setcounter{secnumdepth}{2}
	\phantomsection
	\label{toc}
	\tableofcontents
	
	\vspace{-5mm}
	\section*{Learning objective}\vspace{-5mm}
	Apply minimax regret and Expected Monetary Value (EMV) to select an alternative under uncertainty, and use sensitivity analysis to determine how a subjective probability changes the preferred alternative.
	
	\section*{Criteria assessment}\vspace{-5mm}
	This task sheet assesses Criteria C5 through C7. Each problem assesses the criterion indicated in its title, and each problem is worth 2 points.
	Partial credit may be awarded when the correct procedure is clearly shown but arithmetic errors are made. Answers that do not include the required procedure will not receive credit.
	Failure to write your name on the task sheet or submitting it after the deadline will result in a one-point deduction from each criterion.
	
	A minimum of 4 points is required to pass each criterion.
	
	\section*{Evaluated criteria}
	\begin{enumerate}[label=\textbf{C\arabic*}, start=5]
		\item Applies the Minimax Regret criterion by constructing an opportunity-loss table and selecting the alternative with the smallest maximum regret.
		\item Calculates EMV using given probabilities and selects the alternative with the highest expected payoff.
		\item Uses subjective probabilities and sensitivity analysis to identify threshold probabilities and the preferred alternative over each feasible range.
	\end{enumerate}
	
	\sectionScored{Community Garden Greenhouse}{C5-\sml{2}}{sec:garden-c5}{
		A community garden must choose one of three greenhouse plans. Annual net benefits (in thousands of dollars) depend on the growing season.
		\begin{center}\renewcommand{\arraystretch}{1.25}\begin{tabular}{lccc}\toprule
				Alternative & Excellent ($S_1$) & Average ($S_2$) & Poor ($S_3$)\\ \midrule
				Large greenhouse ($A_1$)&70&30&-10\\
				Medium greenhouse ($A_2$)&55&45&20\\
				Small greenhouse ($A_3$)&35&40&32\\ \bottomrule\end{tabular}\end{center}
		Use the Minimax Regret criterion to determine which plan should be selected. Show the best payoff in each state, the regret table, and each alternative's maximum regret. \score{2}
		
		\subsection*{C5 --- Solution}
		The best payoff in each state is
		\[
		b(S_1)=\max\{70,55,35\}=70,\quad
		b(S_2)=\max\{30,45,40\}=45,\quad
		b(S_3)=\max\{-10,20,32\}=32.
		\]
		Using $r(A_i,S_j)=b(S_j)-P_{ay}(A_i,S_j)$ gives
		\begin{center}\renewcommand{\arraystretch}{1.25}\begin{tabular}{lcccc}\toprule
				& $S_1$ & $S_2$ & $S_3$ & Maximum regret\\ \midrule
				$A_1$&$70-70=0$&$45-30=15$&$32-(-10)=42$&$42$\\
				$A_2$&$70-55=15$&$45-45=0$&$32-20=12$&$15$\\
				$A_3$&$70-35=35$&$45-40=5$&$32-32=0$&$35$\\ \bottomrule\end{tabular}\end{center}
		Thus $\min\{42,15,35\}=15$, so the Minimax Regret criterion selects $\boxed{A_2}$, the medium greenhouse.
		\secInicio[sec:c5-exam-solutions][sec:garden-c5]
	}
	
	\sectionScored{Coastal Ferry Capacity}{C5-\sml{2}}{sec:ferry-c5}{
		A ferry operator is choosing a capacity plan. Seasonal profit (in thousands of dollars) is shown below; no reliable probabilities are available.
		\begin{center}\renewcommand{\arraystretch}{1.25}\begin{tabular}{lccc}\toprule
				Alternative & High demand ($S_1$) & Normal demand ($S_2$) & Low demand ($S_3$)\\ \midrule
				Add two ferries ($A_1$)&100&35&-25\\
				Add one ferry ($A_2$)&75&55&10\\
				Lease peak capacity ($A_3$)&55&48&30\\
				Keep current fleet ($A_4$)&30&30&30\\ \bottomrule\end{tabular}\end{center}
		Determine the decision using Minimax Regret. \score{2}
		
		\subsection*{C5 --- Solution}
		The statewise best payoffs are $b(S_1)=100$, $b(S_2)=55$, and $b(S_3)=30$.
		\begin{center}\renewcommand{\arraystretch}{1.25}\begin{tabular}{lcccc}\toprule
				& $S_1$ & $S_2$ & $S_3$ & Maximum regret\\ \midrule
				$A_1$&$100-100=0$&$55-35=20$&$30-(-25)=55$&$55$\\
				$A_2$&$100-75=25$&$55-55=0$&$30-10=20$&$25$\\
				$A_3$&$100-55=45$&$55-48=7$&$30-30=0$&$45$\\
				$A_4$&$100-30=70$&$55-30=25$&$30-30=0$&$70$\\ \bottomrule\end{tabular}\end{center}
		Since $\min\{55,25,45,70\}=25$, the recommended decision is $\boxed{A_2}$: add one ferry.
		\secInicio[sec:c5-exam-solutions][sec:ferry-c5]
	}
	
	\sectionScored{Mobile Clinic Schedule}{C5-\sml{2}}{sec:clinic-c5}{
		A health organization must select a mobile-clinic schedule. Net community value (in thousands of dollars) depends on attendance.
		\begin{center}\renewcommand{\arraystretch}{1.25}\begin{tabular}{lcccc}\toprule
				Alternative & Very high ($S_1$)&High ($S_2$)&Moderate ($S_3$)&Low ($S_4$)\\ \midrule
				Full schedule ($A_1$)&90&50&20&-20\\
				Flexible schedule ($A_2$)&70&60&45&10\\
				Limited schedule ($A_3$)&45&48&42&35\\ \bottomrule\end{tabular}\end{center}
		Use Minimax Regret to select a schedule. \score{2}
		
		\subsection*{C5 --- Solution}
		The best-payoff vector is
		\[
		(b(S_1),b(S_2),b(S_3),b(S_4))=(90,60,45,35).
		\]
		\begin{center}\renewcommand{\arraystretch}{1.25}\begin{tabular}{lccccc}\toprule
				&$S_1$&$S_2$&$S_3$&$S_4$&Maximum regret\\ \midrule
				$A_1$&0&10&25&55&55\\
				$A_2$&20&0&0&25&25\\
				$A_3$&45&12&3&0&45\\ \bottomrule\end{tabular}\end{center}
		For example, $r(A_1,S_4)=35-(-20)=55$ and $r(A_3,S_2)=60-48=12$. The smallest maximum regret is $\min\{55,25,45\}=25$, so select $\boxed{A_2}$, the flexible schedule.
		\secInicio[sec:c5-exam-solutions][sec:clinic-c5]
	}
	
	\sectionScored{Urban Farm Crop Mix}{C6-\sml{2}}{sec:farm-c6}{
		An urban farm is selecting a crop mix. Profits are in thousands of dollars.
		\begin{center}\renewcommand{\arraystretch}{1.25}\begin{tabular}{lccc}\toprule
				&Strong market ($S_1$)&Stable market ($S_2$)&Weak market ($S_3$)\\
				$P(S_j)$&0.50&0.30&0.20\\ \midrule
				Specialty crops ($A_1$)&80&40&-10\\
				Balanced crops ($A_2$)&60&55&20\\
				Staple crops ($A_3$)&40&45&35\\ \bottomrule\end{tabular}\end{center}
		Calculate every EMV and apply the Bayes decision criterion. \score{2}
		
		\subsection*{C6 --- Solution}
		First, $0.50+0.30+0.20=1$. Then
		\begin{align*}
			EMV(A_1)&=(0.50)(80)+(0.30)(40)+(0.20)(-10)=40+12-2=50,\\
			EMV(A_2)&=(0.50)(60)+(0.30)(55)+(0.20)(20)=30+16.5+4=50.5,\\
			EMV(A_3)&=(0.50)(40)+(0.30)(45)+(0.20)(35)=20+13.5+7=40.5.
		\end{align*}
		The highest EMV is $\max\{50,50.5,40.5\}=50.5$. Therefore, choose $\boxed{A_2}$, the balanced crop mix, with expected profit $\$50{,}500$.
		\secInicio[sec:c6-exam-solutions][sec:farm-c6]
	}
	
	\sectionScored{Museum Exhibition Scale}{C6-\sml{2}}{sec:museum-c6}{
		A museum is deciding the scale of a temporary exhibition. Payoffs are in thousands of dollars.
		\begin{center}\renewcommand{\arraystretch}{1.25}\begin{tabular}{lccc}\toprule
				&High attendance ($S_1$)&Expected attendance ($S_2$)&Low attendance ($S_3$)\\
				$P(S_j)$&0.25&0.45&0.30\\ \midrule
				Major exhibition ($A_1$)&95&35&-20\\
				Standard exhibition ($A_2$)&70&50&15\\
				Compact exhibition ($A_3$)&50&45&32\\
				Permanent collection only ($A_4$)&25&25&25\\ \bottomrule\end{tabular}\end{center}
		Calculate the EMVs and recommend an alternative. \score{2}
		
		\subsection*{C6 --- Solution}
		The probabilities are valid because $0.25+0.45+0.30=1$.
		\begin{align*}
			EMV(A_1)&=(0.25)(95)+(0.45)(35)+(0.30)(-20)=23.75+15.75-6=33.50,\\
			EMV(A_2)&=(0.25)(70)+(0.45)(50)+(0.30)(15)=17.5+22.5+4.5=44.50,\\
			EMV(A_3)&=(0.25)(50)+(0.45)(45)+(0.30)(32)=12.5+20.25+9.6=42.35,\\
			EMV(A_4)&=(0.25)(25)+(0.45)(25)+(0.30)(25)=25.00.
		\end{align*}
		Because $44.50$ is the highest EMV, the Bayes decision is $\boxed{A_2}$, the standard exhibition, with expected payoff $\$44{,}500$.
		\secInicio[sec:c6-exam-solutions][sec:museum-c6]
	}
	
	\sectionScored{Solar Installer Staffing}{C6-\sml{2}}{sec:solar-c6}{
		A solar installer is choosing a staffing plan for next quarter. Payoffs are in thousands of dollars.
		\begin{center}\renewcommand{\arraystretch}{1.25}\begin{tabular}{lcccc}\toprule
				&Surge ($S_1$)&Strong ($S_2$)&Steady ($S_3$)&Slow ($S_4$)\\
				$P(S_j)$&0.15&0.35&0.30&0.20\\ \midrule
				Rapid expansion ($A_1$)&100&60&20&-30\\
				Moderate expansion ($A_2$)&75&65&45&10\\
				Current staff ($A_3$)&50&52&48&40\\ \bottomrule\end{tabular}\end{center}
		Use EMV to determine the preferred staffing plan. \score{2}
		
		\subsection*{C6 --- Solution}
		The probability check is $0.15+0.35+0.30+0.20=1$.
		\begin{align*}
			EMV(A_1)&=(0.15)(100)+(0.35)(60)+(0.30)(20)+(0.20)(-30)=15+21+6-6=36,\\
			EMV(A_2)&=(0.15)(75)+(0.35)(65)+(0.30)(45)+(0.20)(10)=11.25+22.75+13.5+2=49.5,\\
			EMV(A_3)&=(0.15)(50)+(0.35)(52)+(0.30)(48)+(0.20)(40)=7.5+18.2+14.4+8=48.1.
		\end{align*}
		Thus $\max\{36,49.5,48.1\}=49.5$, so choose $\boxed{A_2}$, moderate expansion, with expected payoff $\$49{,}500$.
		\secInicio[sec:c6-exam-solutions][sec:solar-c6]
	}
	
	\sectionScored{Festival Catering Contract}{C7-\sml{2}}{sec:catering-c7}{
		A caterer must choose between a large contract ($A_1$) and a small contract ($A_2$). Payoffs (in thousands of dollars) depend on whether festival attendance is high ($S_1$) or low ($S_2$).
		\begin{center}\renewcommand{\arraystretch}{1.25}\begin{tabular}{lcc}\toprule
				Alternative&High ($S_1$)&Low ($S_2$)\\ \midrule
				Large contract ($A_1$)&80&10\\
				Small contract ($A_2$)&45&35\\ \bottomrule\end{tabular}\end{center}
		Let $P(S_1)=p$ and $P(S_2)=1-p$. Find the threshold probability and state the preferred alternative for every $0\le p\le1$. \score{2}
		
		\subsection*{C7 --- Solution}
		\begin{align*}
			EMV(A_1)&=80p+10(1-p)=10+70p,\\
			EMV(A_2)&=45p+35(1-p)=35+10p.
		\end{align*}
		At the threshold, $10+70p=35+10p$, so $60p=25$ and $p=\frac{5}{12}\approx0.4167$. Since $A_2$ has the higher EMV at $p=0$ and $A_1$ has the higher EMV at $p=1$,
		\[
		\boxed{\begin{cases}
				A_2,&0\le p<\frac{5}{12},\\
				A_1\text{ and }A_2\text{ are tied},&p=\frac{5}{12},\\
				A_1,&\frac{5}{12}<p\le1.
		\end{cases}}
		\]
		\secInicio[sec:c7-exam-solutions][sec:catering-c7]
	}
	
	\sectionScored{Electric-Bike Launch}{C7-\sml{2}}{sec:bike-c7}{
		A retailer is comparing an aggressive launch ($A_1$) with a cautious launch ($A_2$). Payoffs are in thousands of dollars.
		\begin{center}\renewcommand{\arraystretch}{1.25}\begin{tabular}{lccc}\toprule
				Alternative&High demand ($S_1$)&Moderate demand ($S_2$)&Low demand ($S_3$)\\ \midrule
				Aggressive launch ($A_1$)&90&30&-10\\
				Cautious launch ($A_2$)&45&50&25\\ \bottomrule\end{tabular}\end{center}
		Suppose $P(S_1)=p$ and $P(S_2)=\frac{p}{2}$. Determine the feasible range of $p$, the threshold, and the preferred launch over the entire range. \score{2}
		
		\subsection*{C7 --- Solution}
		Because the three state probabilities must sum to one,
		\[
		P(S_3)=1-p-\frac{p}{2}=1-\frac{3p}{2}.
		\]
		Nonnegative probabilities require $p\ge0$ and $1-\frac{3p}{2}\ge0$, so the feasible range is $0\le p\le\frac{2}{3}$.
		\begin{align*}
			EMV(A_1)&=90p+30\left(\frac{p}{2}\right)-10\left(1-\frac{3p}{2}\right)=120p-10,\\
			EMV(A_2)&=45p+50\left(\frac{p}{2}\right)+25\left(1-\frac{3p}{2}\right)=32.5p+25.
		\end{align*}
		Equating the EMVs gives $120p-10=32.5p+25$, hence
		\[
		87.5p=35\quad\Rightarrow\quad p=\frac{2}{5}=0.4.
		\]
		Therefore,
		\[
		\boxed{\begin{cases}
				A_2,&0\le p<\frac{2}{5},\\
				A_1\text{ and }A_2\text{ are tied},&p=\frac{2}{5},\\
				A_1,&\frac{2}{5}<p\le\frac{2}{3}.
		\end{cases}}
		\]
		\secInicio[sec:c7-exam-solutions][sec:bike-c7]
	}
	
	\sectionScored{Rainwater-System Bid}{C7-\sml{2}}{sec:rainwater-c7}{
		An engineering firm must choose between a standard bid ($A_1$) and an innovative bid ($A_2$). Payoffs are in thousands of dollars.
		\begin{center}\renewcommand{\arraystretch}{1.25}\begin{tabular}{lccc}\toprule
				Alternative&Immediate approval ($S_1$)&Revision ($S_2$)&Delayed approval ($S_3$)\\ \midrule
				Standard bid ($A_1$)&20&55&75\\
				Innovative bid ($A_2$)&70&40&30\\ \bottomrule\end{tabular}\end{center}
		Suppose $P(S_1)=2p$ and $P(S_2)=p$. Find the feasible range, threshold probability, and preferred bid for each range of $p$. \score{2}
		
		\subsection*{C7 --- Solution}
		The remaining probability is
		\[
		P(S_3)=1-2p-p=1-3p.
		\]
		The probability conditions $p\ge0$ and $1-3p\ge0$ give the feasible range $0\le p\le\frac{1}{3}$.
		\begin{align*}
			EMV(A_1)&=20(2p)+55p+75(1-3p)=75-130p,\\
			EMV(A_2)&=70(2p)+40p+30(1-3p)=30+90p.
		\end{align*}
		At indifference,
		\[
		75-130p=30+90p
		\quad\Rightarrow\quad 45=220p
		\quad\Rightarrow\quad p=\frac{9}{44}\approx0.2045.
		\]
		Here $A_1$ is preferred below the threshold, while the faster-growing EMV of $A_2$ makes it preferred above the threshold:
		\[
		\boxed{\begin{cases}
				A_1,&0\le p<\frac{9}{44},\\
				A_1\text{ and }A_2\text{ are tied},&p=\frac{9}{44},\\
				A_2,&\frac{9}{44}<p\le\frac{1}{3}.
		\end{cases}}
		\]
		\secInicio[sec:c7-exam-solutions][sec:rainwater-c7]
	}
	
\end{document}
